% Generated draft: requires mathematical and source review.
\documentclass[11pt]{article}
\usepackage[margin=25mm]{geometry}
\usepackage{fontspec}
\setmainfont{Latin Modern Roman}
\usepackage{amsmath,amssymb,mathtools,unicode-math}
\setmathfont{Latin Modern Math}
\usepackage{enumitem,needspace,xurl,hyperref}
\setlength{\parindent}{0pt}
\setlength{\parskip}{5pt}
\setlength{\emergencystretch}{3em}
\newcommand{\N}{\mathbb{N}}
\newcommand{\Z}{\mathbb{Z}}
\newcommand{\Q}{\mathbb{Q}}
\newcommand{\R}{\mathbb{R}}
\newcommand{\C}{\mathbb{C}}
\newcommand{\F}{\mathbb{F}}
\newcommand{\A}{\mathbb{A}}
\newcommand{\PP}{\mathbb{P}}
\newcommand{\E}{\mathbb{E}}
\newcommand{\eps}{\varepsilon}
\newcommand{\dd}{\,\mathrm{d}}
\begin{document}
\section*{UnsolvedMath ID: 30003359}
\textbf{Source problem code:} catalogue entry
\subsection{Definitions and mathematical statement}
Let \((M,g)\) be a smooth Riemannian manifold. Here \emph{closed} means compact and without boundary, and \emph{Ricci-flat} means that its Ricci curvature tensor satisfies \(\operatorname{Ric}_g=0\). The supplied statement imposes no dimension restriction and does not specify a spin structure.

To express the spinor question, suppose additionally that \(M\) admits a spin structure \(\sigma\). Let \(\Sigma_{g,\sigma}M\) denote its complex spinor bundle, let \(\nabla^{g,\sigma}\) be the spin connection induced by the Levi--Civita connection, and let \(D_{g,\sigma}\) be the associated Dirac operator. In a local orthonormal frame \(e_1,\ldots,e_n\),
\[
D_{g,\sigma}\psi=\sum_{j=1}^{n}e_j\mathbin{\cdot}\nabla^{g,\sigma}_{e_j}\psi,
\]
where \(\mathbin{\cdot}\) denotes Clifford multiplication. A spinor is globally parallel if \(\nabla^{g,\sigma}_X\psi=0\) for every smooth vector field \(X\) on \(M\), and is harmonic if \(D_{g,\sigma}\psi=0\).

The supplied question is: if \((M,g)\) is closed and Ricci-flat, does it admit a globally parallel harmonic spinor? Under the additional interpretation that a spin structure \(\sigma\) has been specified and that the requested spinor must be nonzero, the existence condition is
\[
\exists\,\psi\in\Gamma^{\infty}(\Sigma_{g,\sigma}M),\qquad
\psi\not\equiv0,\qquad
\nabla^{g,\sigma}\psi=0,\qquad
D_{g,\sigma}\psi=0.
\]
The harmonic condition follows directly from the parallel condition by the displayed definition of \(D_{g,\sigma}\); it is retained to match the source wording.

\textbf{Unresolved formulation qualifications.} The original and cleaned statements do not assume that \(M\) is spin, whereas the supplied literature assessment describes a question about closed Ricci-flat \emph{spin} metrics. The record does not specify whether spinness is an intended hypothesis or part of the requested conclusion. Even under a spin hypothesis, it does not specify whether the conclusion concerns a prescribed spin structure, every spin structure, or the existence of some spin structure admitting such a spinor. These are distinct quantifier choices, so the displayed condition is a conditional formulation, not a resolution of the missing quantifier. The source also omits the word ``nonzero''; that requirement is an explicit interpretation needed to exclude the zero spinor. No passage to a covering manifold is stated or inserted.
\subsection{Short English statement}
The question asks whether every closed Ricci-flat manifold carries a nonzero globally parallel harmonic spinor. A precise version requires clarification of the spin hypothesis and the quantifier over spin structures. The nonzero requirement is implicit rather than explicit in the supplied wording. The registry labels the problem partially solved, but that label and its literature assessment are supplied curation requiring independent review.
\subsection{Sources}
\begin{itemize}
\item Supplied local UnsolvedMath registry record, permanent ID \(30003359\), source problem code \texttt{catalogue entry}, titled \emph{Parallel Harmonic Spinors on Closed Ricci-Flat Manifolds}, in the collection \emph{Oberwolfach Reports Open Problems}. No individual proposer is supplied. The proposed year is recorded as 2017.
\item Supplied source citation: \emph{Mini-Workshop: Spaces and Moduli Spaces of Riemannian Metrics (2018)}, Oberwolfach Reports, pp.\ 133--166, DOI \texttt{10.4171/owr/2017/3}. \url{https://doi.org/10.4171/owr/2017/3}. The background describes Oberwolfach Reports as documenting mathematical workshops and their open-problem discussions. The literature-source label for this link is ``Oberwolfach report on Ricci-flat metrics.'' The citation's parenthetical year 2018 differs from the recorded proposed year and the year appearing in the DOI; this discrepancy remains unresolved.
\item Supplied literature reference: \emph{Ricci-flat manifolds, parallel spinors, and the Rosenberg index}. \url{https://doi.org/10.3842/SIGMA.2025.072}.
\item Supplied literature reference: \emph{Holonomy rigidity for Ricci-flat metrics}. \url{https://arxiv.org/abs/1512.07390}.
\item Literature assessment dated exactly \texttt{2026-08-21}: the supplied curation labels the problem ``partially solved'' and states that the general converse for closed Ricci-flat spin metrics remains open, while a 2025 theorem establishes it under a nonzero Rosenberg-index hypothesis and yields the expected parallel-spinor conclusion in several major classes. This is supplied dated curation, not independently verified here. The record does not specify the full theorem hypotheses, the relevant spin-structure quantifier, or the classes mentioned. Its background explicitly calls for independent verification of time-sensitive status and cited claims.
\item The registry assigns the statement status \texttt{corrected\_verified}, with the supplied explanation that notation, typography, and sentence structure were repaired without changing mathematical content. This editorial status does not establish that the missing mathematical qualifications have been resolved.
\end{itemize}
\end{document}
